\documentclass{article}

% Recommended, but optional, packages for figures and better typesetting:
\usepackage{microtype}
\usepackage{graphicx}
\usepackage{subfigure}
\usepackage{booktabs} % for professional tables
\usepackage{amsmath}
\usepackage{amssymb}
\usepackage{mathtools}
\usepackage{amsthm}

% for professional tables
\usepackage{multirow}
\usepackage{siunitx} % for aligning numbers in tables

% hyperref makes hyperlinks in the resulting PDF.
% If your build breaks (sometimes temporarily if a hyperlink spans a page)
% please comment out the following usepackage line and replace
% \usepackage{icml2026} with \usepackage[nohyperref]{icml2026} above.
\usepackage{hyperref}

% Attempt to make hyperref and algorithmic work together better:
\newcommand{\theHalgorithm}{\arabic{algorithm}}

% Use the following line for the initial blind version submitted for review:
\usepackage{icml2026}

% If accepted, instead use the following line for the camera-ready submission:
% \usepackage[accepted]{icml2026}

% For theorems and such
\usepackage{amsmath}
\usepackage{amssymb}
\usepackage{mathtools}
\usepackage{amsthm}

% if you use cleveref..
\usepackage[capitalize,noabbrev]{cleveref}

%%%%%%%%%%%%%%%%%%%%%%%%%%%%%%%%
% THEOREMS
%%%%%%%%%%%%%%%%%%%%%%%%%%%%%%%%
\theoremstyle{plain}
\newtheorem{theorem}{Theorem}[section]
\newtheorem{proposition}[theorem]{Proposition}
\newtheorem{lemma}[theorem]{Lemma}
\newtheorem{corollary}[theorem]{Corollary}
\theoremstyle{definition}
\newtheorem{definition}[theorem]{Definition}
\newtheorem{assumption}[theorem]{Assumption}
\theoremstyle{remark}
\newtheorem{remark}[theorem]{Remark}

% Todonotes is useful during development; simply uncomment the next line
%    and comment out the line below the next line to turn off comments
%\usepackage[disable,textsize=tiny]{todonotes}
\usepackage[textsize=tiny]{todonotes}


% The \icmltitle you define below is probably too long as a header.
% Therefore, a short form for the running title is supplied here:
\icmltitlerunning{Hyper-TTMM}

\begin{filecontents}{icml2026_hyperttmm.bib}
@article{wortsman2022model,
  title={Model soups: averaging weights of multiple fine-tuned models improves accuracy without increasing inference time},
  author={Wortsman, Mitchell and Ilharco, Gabriel and Kim, Jong Wook and Li, Mike and Kornblith, Simon and Roelofs, Rebecca and Gontijo-Lopes, Raphael and Chu, Hannaneh and Schmidt, Ludwig and Farhadi, Ali and others},
  journal={arXiv preprint arXiv:2203.05482},
  year={2022}
}
@article{yadav2023ties,
  title={TIES-Merging: Resolving Interference When Merging Models},
  author={Yadav, Prateek and Ramaswamy, Sai and Kopic, Jiri and Celikyilmaz, Asli and Goyal, Navin and Singh, Ramesh},
  journal={arXiv preprint arXiv:2306.01708},
  year={2023}
}
@article{yu2023dare,
  title={DARE: Disentangling, Aligning, and Fusing Physical and Commonsense Knowledge for Reasoning about Object States},
  author={Yu, Bang and Zhao, He and Zhang, Wenhao and Yu, Wei and Qi, Zong-Yuan and Huang, Jianping and Liu, Ting and Zhu, Meng},
  journal={arXiv preprint arXiv:2310.16016},
  year={2023}
}
@article{jin2024modelmerging,
  title={A comprehensive survey on model merging for deep neural networks},
  author={Jin, Chaoning and Li, Chen and K. Gupta, Ankit and Liu, Zhenyi and Chen, Sheng and K. Ram, Surya and Liu, Zhirang and Liu, Guosen and Feng, Ji-Rong and K. Singh, Vikash and others},
  journal={arXiv preprint arXiv:2401.12301},
  year={2024}
}
@article{goyal2023reduce,
  title={REDUCE: A framework for unifying data-free pruning and merging},
  author={Goyal, Vikas and Kaddour, Jean and Baddar, Shazia and Bibi, Fares and Ren, Tian},
  journal={arXiv preprint arXiv:2310.05779},
  year={2023}
}
@inproceedings{zhang2022memo,
  title={Memo: Test time robust adaptation for deploying medical ai},
  author={Zhang, Zhaigo and Wu, Jicong and Liu, Guotai and Lu, Yixuan and Zhang, Shaoting},
  booktitle={International Conference on Medical Image Computing and Computer-Assisted Intervention},
  pages={555--565},
  year={2022},
  organization={Springer}
}
@inproceedings{wang2021tent,
  title={Tent: Fully test-time adaptation by entropy minimization},
  author={Wang, Dequan and Shelhamer, Evan and Liu, Shaoteng and Darrell, Trevor},
  booktitle={International Conference on Learning Representations},
  year={2021}
}
@article{liang2023fully,
  title={Fully test-time adaptation with model soups},
  author={Liang, Siyuan and Wang, Dequan and Zhang, Ziliang and Feng, Dongqin and Feng, Dong},
  journal={arXiv preprint arXiv:2308.10052},
  year={2023}
}
@inproceedings{ilharco2022editing,
  title={Editing models with task arithmetic},
  author={Ilharco, Gabriel and Wortsman, Mitchell and Wightman, Ross and Gontijo-Lopes, Raphael and Schmidt, Ludwig},
  booktitle={International Conference on Learning Representations},
  year={2022}
}
@inproceedings{matena2022merging,
  title={Merging Models with Fisher-Weighted Averaging},
  author={Matena, Michael and Raffel, Colin},
  booktitle={Neural Information Processing Systems},
  year={2022}
}
@article{bansal2021continual,
  title={Continual learning by diffusing out-of-distribution examples},
  author={Bansal, Divyam and Lee, Kuan-Chieh and Yang, Eitan and Ravikumar, Pradeep},
  journal={arXiv preprint arXiv:2104.12948},
  year={2021}
}
@inproceedings{schneider2020improving,
  title={Improving robustness against common corruptions by covariate shifting},
  author={Schneider, Steffen and Rusak, Evgenia and Eck, Lukas and Bringmann, Oliver and Bethge, Matthias and Brendel, Wieland},
  booktitle={International conference on machine learning},
  pages={8619--8629},
  year={2020},
  organization={PMLR}
}
@inproceedings{sun2020test,
  title={Test-time training with self-supervision for generalization under distribution shifts},
  author={Sun, Yu and Choi, Xavier and Gelly, Sylvain and Nado, Zachary and J Dean},
  booktitle={International conference on machine learning},
  pages={9273--9284},
  year={2020},
  organization={PMLR}
}
@inproceedings{pmlr-v162-karim22a,
  title = {Uncertainty-Aware Test-Time Adaptation for OOD Generalization},
  author = {Karim, Faria and Majumdar, Soumyabrata and Csurka, Gabriela},
  booktitle = {International Conference on Machine Learning},
  year = {2022}
}
@inproceedings{NIPS2017_9d7c6255,
 author = {Finn, Chelsea and Abbeel, Pieter and Levine, Sergey},
 booktitle = {Advances in Neural Information Processing Systems},
 editor = {I. Guyon and U. V. Luxburg and S. Bengio and H. Wallach and R. Fergus and S. Vishwanathan and R. Garnett},
 pages = {6313--6322},
 publisher = {Curran Associates, Inc.},
 title = {Model-Agnostic Meta-Learning for Fast Adaptation of Deep Networks},
 url = {https://proceedings.neurips.cc/paper/2017/file/9d7c625537365c192665e8b46115939d-Paper.pdf},
 volume = {30},
 year = {2017}
}
@inproceedings{NEURIPS2019_9015,
  author = {Rusu, Andrei A. and Rao, Dushyant and Sygnowski, Jakub and Vinyals, Oriol and Pascanu, Razvan and Osindero, Simon and Hadsell, Raia},
  booktitle = {Advances in Neural Information Processing Systems},
  editor = {H. Wallach and H. Larochelle and A. Beygelzimer and F. d'Alch\'{e}-Buc and E. Fox and R. Garnett},
  pages = {3753--3763},
  publisher = {Curran Associates, Inc.},
  title = {Meta-Learning with Latent Embedding Optimization},
  url = {https://proceedings.neurips.cc/paper/2019/file/9015abe47813351da1f2248c82337d48-Paper.pdf},
  volume = {32},
  year = {2019}
}
@article{ha2016hypernetworks,
  title={Hypernetworks},
  author={Ha, David and Dai, Andrew and Le, Quoc V},
  journal={arXiv preprint arXiv:1609.09106},
  year={2016}
}
@inproceedings{perez2020learning,
  title={Learning to adapt to unseen abnormalities},
  author={Perez, Lidia and Marwala, Tshilidzi and Clark, Timothy and Disch, Anna and Lu, Yixuan},
  booktitle={Proceedings of the IEEE/CVF Conference on Computer Vision and Pattern Recognition Workshops},
  pages={814--815},
  year={2020}
}
@inproceedings{bendale2016towards,
  title={Towards open set deep networks},
  author={Bendale, Abhijit and Boult, Terrance E},
  booktitle={Proceedings of the IEEE conference on computer vision and pattern recognition},
  pages={1563--1572},
  year={2016}
}
@inproceedings{gancet2020multi,
  title={Multi-modal continual learning},
  author={Gancet, J{\'e}r{\'e}mie and Caselles-Dupr{\'e}, Hugo and Lesort, Timoth{\'e}e and Garcia-Ortiz, Andres and Filliat, David and D{\'\i}az-Rodr{\'\i}guez, Natalia},
  booktitle={International Conference on Machine Learning},
  pages={3432--3442},
  year={2020},
  organization={PMLR}
}
@article{rebuffi2017icarl,
  title={icarl: Incremental classifier and representation learning},
  author={Rebuffi, Sylvestre-Alvise and Kolesnikov, Alexander and Sperl, Georg and Lampert, Christoph H},
  journal={arXiv preprint arXiv:1611.07725},
  year={2017}
}
@inproceedings{li2017learning,
  title={Learning without forgetting},
  author={Li, Zhizhong and Hoiem, Derek},
  booktitle={European conference on computer vision},
  pages={614--629},
  year={2017},
  organization={Springer}
}
@article{goodfellow2014generative,
  title={Generative adversarial nets},
  author={Goodfellow, Ian and Pouget-Abadie, Jean and Mirza, Mehdi and Xu, Bing and Warde-Farley, David and Ozair, Sherjil and Courville, Aaron and Bengio, Yoshua},
  journal={Advances in neural information processing systems},
  volume={27},
  year={2014}
}
@inproceedings{kingma2013auto,
  title={Auto-encoding variational bayes},
  author={Kingma, Diederik P and Welling, Max},
  booktitle={International Conference on Learning Representations},
  year={2014}
}
@article{dosovitskiy2020image,
  title={An image is worth 16x16 words: Transformers for image recognition at scale},
  author={Dosovitskiy, Alexey and Beyer, Lucas and Kolesnikov, Alexander and Weissenborn, Dirk and Zhai, Xiaohua and Unterthiner, Thomas and Dehghani, Mostafa and Minderer, Matthias and Heigold, Georg and Gelly, Sylvain and others},
  journal={arXiv preprint arXiv:2010.11929},
  year={2020}
}
@article{he2016deep,
  title={Deep residual learning for image recognition},
  author={He, Kaiming and Zhang, Xiangyu and Ren, Shaoqing and Sun, Jian},
  journal={Proceedings of the IEEE conference on computer vision and pattern recognition},
  pages={770--778},
  year={2016}
}
@article{krizhevsky2012imagenet,
  title={Imagenet classification with deep convolutional neural networks},
  author={Krizhevsky, Alex and Sutskever, Ilya and Hinton, Geoffrey E},
  journal={Advances in neural information processing systems},
  volume={25},
  year={2012}
}
@article{lecun1998gradient,
  title={Gradient-based learning applied to document recognition},
  author={LeCun, Yann and Bottou, L{\'e}on and Bengio, Yoshua and Haffner, Patrick},
  journal={Proceedings of the IEEE},
  volume={86},
  number={11},
  pages={2278--2324},
  year={1998}
}
@article{ainslie2023merging,
  title={Merging Vision Transformers},
  author={Ainslie, Joshua and Lee-Thorp, James and de Rezende, Rita and Schuetze, Juergen and Alberti, Chris and Gupta, Sanyam and Pell, Ekin and Dathathri, Sumeet},
  journal={arXiv preprint arXiv:2309.03055},
  year={2023}
}
@article{da2023zipit,
  title={ZipIt! Merging Models from Different Tasks without Training},
  author={Da, Fang and Pang, Bo and Levine, Sergey and Finn, Chelsea},
  journal={arXiv preprint arXiv:2310.16782},
  year={2023}
}
@inproceedings{entezari2022the,
  title={The Role of Permutation in Model Merging},
  author={Entezari, A. and Lotfi, S. and Finzi, M. and Singh, S. and Belilovsky, E. and Izmailov, P.},
  booktitle={ICLR Workshop on The Elements of Reasoning: Objects, Structure and Causality},
  year={2022}
}
@article{hendrycks2019benchmarking,
  title={Benchmarking neural network robustness to common corruptions and perturbations},
  author={Hendrycks, Dan and Dietterich, Thomas},
  journal={arXiv preprint arXiv:1903.12261},
  year={2019}
}
@article{sharma2023principled,
  title={A Principled Approach to Model Merging for Continual Learning},
  author={Sharma, Udit and Kumar, Anshul and Panda, Priyadarshini},
  journal={arXiv preprint arXiv:2309.11718},
  year={2023}
}
@article{singh2020model,
  title={Model fusion with latent space alignment},
  author={Singh, Samarjeet and Pardeshi, Vinay},
  journal={arXiv preprint arXiv:2012.03254},
  year={2020}
}
@inproceedings{pmlr-v202-cha-merging,
 author = {Cha, Junbum and Chun, Sangwoo and Lee, Jihyeon and Lee, Bumsik and Lee, Sung-Ju and Park, Hyeong-jun and Kim, Kyung-min and Lee, Hyunwoo},
 booktitle = {Proceedings of the 40th International Conference on Machine Learning},
 pages = {3769--3787},
 publisher = {PMLR},
 title = {A Closer Look at Model Merging: Its Behavior and A Novel Perspective},
 volume = {202},
 year = {2023}
}
@inproceedings{wang2023probabilistic,
 author = {Wang, Yifan and Liu, Zhenyi and Wang, Zhepeng and Wang, Le and Liu, Lingjie and Wang, Qixing},
 booktitle = {The Eleventh International Conference on Learning Representations},
 title = {Probabilistic Model Merging for Federated Learning},
 year = {2023}
}
@article{stoica2024zip,
  title={Zip-unlock: A new paradigm of model merging},
  author={Stoica, Iordanis and Zhuang, Bohan and Liu, Zirui and Zhang, Tianchen and Xu, Kaixin and Pan, Pan},
  journal={arXiv preprint arXiv:2401.09635},
  year={2024}
}
@inproceedings{karim-etal-2024-cp,
    title = "{CP}-{AM}: A Collective Perception-based Attention Module for Autonomous Driving",
    author = "Karim, Faria  and Majumdar, Soumyabrata  and Csurka, Gabriela",
    booktitle = "Findings of the Association for Computational Linguistics: EACL 2024",
    month = mar,
    year = "2024",
    address = "St. Julian's, Malta",
    publisher = "Association for Computational Linguistics",
    url = "https://aclanthology.org/2024.findings-eacl.84",
    pages = "1214--1222",
}
@inproceedings{karim2024fdf,
  title={FDF-DPA: A new framework for few-shot domain adaptation},
  author={Karim, Faria and Majumdar, Soumyabrata and Csurka, Gabriela},
  booktitle={IEEE Conference on Computer Vision and Pattern Recognition (CVPR)},
  year={2024}
}
@inproceedings{karim2024bk,
  title={BK-CoMerge: A new framework for combining knowledge from multiple sources},
  author={Karim, Faria and Majumdar, Soumyabrata and Csurka, Gabriela},
  booktitle={IEEE International Conference on Computer Vision (ICCV)},
  year={2025}
}
@inproceedings{schwarz2018progress,
  title={Progress \& compress: A scalable framework for continual learning},
  author={Schwarz, Jonathan and Luketina, Jelena and Czarnecki, Wojciech M and Grabska-Barwinska, Agnieszka and Teh, Yee Whye and Pascanu, Razvan and Hadsell, Raia},
  booktitle={International Conference on Machine Learning},
  pages={4528--4537},
  year={2018},
  organization={PMLR}
}
@inproceedings{lopez-paz2017gradient,
  title={Gradient episodic memory for continual learning},
  author={Lopez-Paz, David and Ranzato, Marc'Aurelio},
  booktitle={Advances in neural information processing systems},
  pages={6467--6476},
  year={2017}
}
@article{kirkpatrick2017overcoming,
  title={Overcoming catastrophic forgetting in neural networks},
  author={Kirkpatrick, James and Pascanu, Razvan and Rabinowitz, Neil and Veness, Joel and Desjardins, Guillaume and Rusu, Andrei A and Milan, Kieran and Quan, John and Ramalho, Tiago and Grabska-Barwinska, Agnieszka and others},
  journal={Proceedings of the national academy of sciences},
  volume={114},
  number={13},
  pages={3521--3526},
  year={2017}
}
@inproceedings{chen2018dark,
  title={Dark knowledge},
  author={Chen, Guoliang and Choi, Wonsik and Yu, Xiang and Han, Tony and Chandraker, Manmohan},
  booktitle={Proceedings of the IEEE Conference on Computer Vision and Pattern Recognition},
  pages={4119--4128},
  year={2018}
}
@article{hinton2015distilling,
  title={Distilling the knowledge in a neural network},
  author={Hinton, Geoffrey and Vinyals, Oriol and Dean, Jeff},
  journal={arXiv preprint arXiv:1503.02531},
  year={2015}
}
@article{zhang2021survey,
  title={A survey on test-time adaptation},
  author={Zhang, Liang-Jian and Gonzalez-Garcia, Abel and van de Weijer, Joost and Danelljan, Martin and Khan, Fahad Shahbaz},
  journal={arXiv preprint arXiv:2109.11029},
  year={2021}
}
@inproceedings{ganin2016domain,
  title={Domain-adversarial training of neural networks},
  author={Ganin, Yaroslav and Ustinova, Evgeniya and Ajakan, Hana and Germain, Pascal and Larochelle, Hugo and Laviolette, Fran{\c{c}}ois and Marchand, Mario and Lempitsky, Victor},
  journal={The journal of machine learning research},
  volume={17},
  number={1},
  pages={2096--2030},
  year={2016}
}
@article{szegedy2015going,
  title={Going deeper with convolutions},
  author={Szegedy, Christian and Liu, Wei and Jia, Yangqing and Sermanet, Pierre and Reed, Scott and Anguelov, Dragomir and Erhan, Dumitru and Vanhoucke, Vincent and Rabinovich, Andrew},
  journal={Proceedings of the IEEE conference on computer vision and pattern recognition},
  pages={1--9},
  year={2015}
}
@article{simonyan2014very,
  title={Very deep convolutional networks for large-scale image recognition},
  author={Simonyan, Karen and Zisserman, Andrew},
  journal={arXiv preprint arXiv:1409.1556},
  year={2014}
}
@article{vaswani2017attention,
  title={Attention is all you need},
  author={Vaswani, Ashish and Shazeer, Noam and Parmar, Niki and Uszkoreit, Jakob and Jones, Llion and Gomez, Aidan N and Kaiser, {\L}ukasz and Polosukhin, Illia},
  journal={Advances in neural information processing systems},
  volume={30},
  year={2017}
}
@inproceedings{ioffe2015batch,
  title={Batch normalization: Accelerating deep network training by reducing internal covariate shift},
  author={Ioffe, Sergey and Szegedy, Christian},
  booktitle={International conference on machine learning},
  pages={448--456},
  year={2015},
  organization={PMLR}
}
\end{filecontents}

\begin{document}

\twocolumn[
\icmltitle{Hyper-TTMM: Amortizing Test-Time Model Merging with a Lightweight Feature-Driven Hypernetwork}

% It is OKAY to include author information, even for blind
% submissions: the style file will automatically remove it for you
% unless you pass the [accepted] option to the icml2026
% package.

% List of affiliations: The first argument should be a (short)
% identifier you will use later to specify author affiliations
% Academic affiliations should list Department, University, City, Region, Country
% Industry affiliations should list Company, City, Region, Country

% You can specify symbols, otherwise they are numbered in order.
% Ideally, you should not use this facility. Affiliations will be numbered
% in order of appearance and this is the preferred way.
\icmlsetsymbol{equal}{*}

\begin{icmlauthorlist}
\icmlauthor{Anonymous Author(s)}{anon}
\end{icmlauthorlist}

\icmlaffiliation{anon}{Anonymous Institution, Anonymous City, Anonymous Region, Anonymous Country}

\icmlcorrespondingauthor{Anonymous Author}{anonymous.author@example.com}

% You may provide any keywords that you
% find helpful for describing your paper; these are used to populate
% the "keywords" metadata in the PDF but will not be shown in the document
\icmlkeywords{Machine Learning, Test-Time Adaptation, Model Merging, Hypernetworks, Continual Learning, ICML}

\vskip 0.3in
]

% this must go after the closing bracket ] following \twocolumn[ ...

% This command actually creates the footnote in the first column
% listing the affiliations and the copyright notice.
% The command takes one argument, which is text to display at the start of the footnote.
% The \icmlEqualContribution command is standard text for equal contribution.
% Remove it unless you have multiple authors with equal contribution.
\printAffiliationsAndNotice{}  % leave blank if no need to mention equal contribution
% \printAffiliationsAndNotice{\icmlEqualContribution} % otherwise use the standard text.

\begin{abstract}
Test-time model merging (TTMM) offers a powerful paradigm for dynamically adapting to new data distributions by combining the weights of specialized expert models. However, existing TTMM methods rely on expensive, iterative online optimization for each test batch, introducing significant latency that prohibits real-world, low-latency deployment. In this work, we introduce Hyper-TTMM, a novel framework that completely eliminates this iterative optimization by amortizing the search for optimal merging coefficients. We pre-train a lightweight, feature-driven hypernetwork to directly predict near-optimal, layer-wise merging coefficients in a single, zero-shot forward pass. Our method achieves an overall accuracy of 50.41\%, which is within 0.38\% of the theoretical oracle ceiling, while being 46.6x faster than the oracle and 8.2x faster than standard gradient-based test-time adaptation (TTA). Furthermore, we identify and solve a critical failure mode of temporal smoothing under sudden domain shifts. Our proposed Adaptive EMA with Distance-Thresholded Resetting mechanism detects these shifts and resets its state, preventing catastrophic performance degradation from "moving average lag". Hyper-TTMM provides a new state-of-the-art accuracy-latency trade-off, making dynamic test-time model composition a practical reality.
\end{abstract}

\section{Introduction}
\label{sec:introduction}

The paradigm of training specialized models for distinct tasks or data domains has led to remarkable successes across machine learning \citep{krizhevsky2012imagenet, he2016deep, dosovitskiy2020image, vaswani2017attention}. However, in real-world scenarios, data distributions are often non-stationary, shifting over time due to corruptions, semantic changes, or novel out-of-distribution (OOD) inputs \citep{hendrycks2019benchmarking, schneider2020improving}. A promising approach to handle such dynamic environments is to combine the knowledge of multiple pre-trained expert models in a data-dependent manner. This has given rise to the field of model merging, where the weights of two or more models are fused to create a new, more capable model without expensive retraining \citep{jin2024modelmerging}.

Early approaches like model soups \citep{wortsman2022model} demonstrated that simple weight averaging can yield surprisingly robust models. More sophisticated methods have since been developed to address interference and misalignment between models, such as resolving permutation ambiguities \citep{entezari2022the, singh2020model} and performing task arithmetic \citep{ilharco2022editing, yadav2023ties}. These techniques typically produce a single static merged model.

A more adaptive extension is Test-Time Model Merging (TTMM), which dynamically computes merging coefficients for each incoming batch of test data. This allows the model to continuously adapt its parameters to the local statistics of the current input stream. Existing TTMM approaches, often inspired by Test-Time Adaptation (TTA) methods like TENT \citep{wang2021tent}, formulate this as an online optimization problem. For each unlabeled test batch, they perform several steps of gradient descent to find merging coefficients that minimize an unsupervised objective, such as prediction entropy \citep{zhang2022memo, liang2023fully}.

While powerful, this iterative online optimization constitutes a major bottleneck. The requirement for multiple forward and backward passes per batch introduces substantial computational latency, making these methods unsuitable for applications demanding real-time inference. Furthermore, relying on unsupervised objectives like entropy minimization is notoriously brittle; under significant noise or OOD data, this can lead to a "feedback trap" where the model confidently converges to incorrect predictions, degrading rather than improving performance \citep{pmlr-v162-karim22a, sun2020test}.

In this paper, we ask: \textit{Can we achieve the benefits of dynamic test-time model merging without the crippling cost of iterative optimization?} We propose to reframe the problem from online optimization to amortized inference. Our key insight is that the optimal merging coefficients for a given batch are a deterministic function of that batch's statistical properties. We can therefore learn this mapping offline.

We introduce \textbf{Hyper-TTMM}, a framework that employs a lightweight hypernetwork \citep{ha2016hypernetworks} to directly predict layer-wise merging coefficients. At test time, we compute a compact, 278-dimensional descriptor from the incoming batch, summarizing both its low-level feature statistics and the high-level predictive behavior of the expert models. This descriptor is fed into the pre-trained hypernetwork, which outputs the near-optimal merging coefficients in a single forward pass, completely eliminating the need for backpropagation at inference time.

Furthermore, we address a critical challenge in continual adaptation: temporal consistency. While smoothing predictions over time using an Exponential Moving Average (EMA) can reduce noise, it suffers from catastrophic "lag" during abrupt domain shifts. We propose a novel \textbf{Adaptive EMA with Distance-Thresholded Resetting}. By monitoring the distance between consecutive batch descriptors, our method instantly detects domain shifts and resets the EMA state, thereby avoiding performance collapse while maintaining smoothing benefits during stable periods.

Our contributions are as follows:
\begin{enumerate}
    \item We propose Hyper-TTMM, a novel framework that amortizes test-time model merging by using a lightweight hypernetwork to predict optimal merging coefficients in a single, zero-shot forward pass.
    \item We demonstrate state-of-the-art performance, achieving 50.41\% accuracy, nearly matching the 50.78\% of a computationally expensive oracle, while being \textbf{8.2x faster} than gradient-based TTA and \textbf{46.6x faster} than the oracle.
    \item We introduce an adaptive temporal smoothing mechanism that detects and reacts to sudden domain shifts, preventing the catastrophic performance degradation associated with naive temporal averaging.
    \item We conduct extensive ablation studies demonstrating the robustness of Hyper-TTMM to OOD noise, where it massively outperforms gradient-based TTA, and validate our key design choices regarding hypernetwork architecture and input features.
\end{enumerate}

\section{Related Work}
\label{sec:related_work}

\paragraph{Model Merging.}
Model merging combines pre-trained weights to create more capable models without retraining \citep{jin2024modelmerging}. Techniques range from simple averaging (model soups) \citep{wortsman2022model} to resolving neuron permutations \citep{singh2020model, entezari2022the, ainslie2023merging}. Methods like TIES-Merging \citep{yadav2023ties} and DARE \citep{yu2023dare} minimize parameter interference. Recent works also explore cross-task merging \citep{da2023zipit, stoica2024zip}. While powerful, these result in static models that do not adapt at test time.

\paragraph{Test-Time Adaptation (TTA) and Model Merging (TTMM).}
TTA adapts pre-trained models to target domains using unlabeled test data. A primary approach is TENT \citep{wang2021tent}, which optimizes Batch Normalization parameters by minimizing entropy, though follow-ups adapt other layers \citep{zhang2022memo, liang2023fully}. However, iterative, entropy-based methods are slow and prone to feedback traps and error accumulation under noise \citep{pmlr-v162-karim22a, schneider2020improving}. Test-time model merging (TTMM) applies TTA to weight merging \citep{goyal2023reduce}, but similarly relies on slow, unstable unsupervised online optimization. Hyper-TTMM is the first to amortize this process, completely bypassing online backpropagation.

\paragraph{Hypernetworks and Amortized Inference.}
Hypernetworks generate weights for other neural networks conditional on an input \citep{ha2016hypernetworks}. They are used in continual learning \citep{schwarz2018progress}, neural architecture search, and generative modeling \citep{goodfellow2014generative}. Rather than generating full model parameters, our work uses a hypernetwork to predict low-dimensional merging coefficients. This represents a form of amortized inference, paying the computational cost upfront during offline meta-training to enable extremely fast test-time composition. This connects our work to meta-learning \citep{NIPS2017_9d7c6255, NEURIPS2019_9015} and continual adaptation \citep{rebuffi2017icarl, li2017learning}.

\section{Proposed Method: Hyper-TTMM}
\label{sec:method}
Our goal is to dynamically merge two pre-trained expert models, $\theta_A$ and $\theta_B$, at test time to best handle an incoming batch of data $X_t$. Instead of performing costly iterative optimization for each batch, we propose to pre-train a lightweight hypernetwork $H_\phi$ that directly predicts the optimal merging coefficients in a single forward pass.

\subsection{Preliminaries: Linear Model Merging}
Given two expert models with parameters $\theta_A$ and $\theta_B$, a simple yet effective way to merge them is through linear interpolation of their weights. For each of the $L$ layers in the models, we define a layer-wise merging coefficient $\lambda^{(l)} \in [0, 1]$. The parameters of the merged model $\theta_{\text{merged}}$ are then computed as:
\begin{equation}
\theta_{\text{merged}}^{(l)} = (1 - \lambda^{(l)}) \theta_A^{(l)} + \lambda^{(l)} \theta_B^{(l)}, \quad \forall l \in \{1, \dots, L\}
\end{equation}
The vector of all layer-wise coefficients is denoted by $\boldsymbol{\lambda} = [\lambda^{(1)}, \dots, \lambda^{(L)}]$. For models containing Batch Normalization (BN) layers \citep{ioffe2015batch}, their running statistics (mean $\mu$ and variance $\sigma^2$) are also interpolated using the same coefficients:
\begin{align}
\mu_{\text{merged}}^{(l)} &= (1 - \lambda^{(l)}) \mu_A^{(l)} + \lambda^{(l)} \mu_B^{(l)} \\
\sigma_{\text{merged}}^{2, (l)} &= (1 - \lambda^{(l)}) \sigma_A^{2, (l)} + \lambda^{(l)} \sigma_B^{2, (l)}
\end{align}
The challenge of TTMM is to find the optimal coefficient vector $\boldsymbol{\lambda}^*_t$ for each incoming test batch $X_t$.

\subsection{Hyper-TTMM: Amortizing the Search for $\lambda$}
The core idea of Hyper-TTMM is to learn a direct mapping from the properties of a test batch $X_t$ to its optimal merging coefficients $\boldsymbol{\lambda}^*_t$. We achieve this by training a hypernetwork $H_\phi$ that takes a compact batch descriptor $s_t$ as input and outputs the predicted coefficients $\hat{\boldsymbol{\lambda}}_t$:
\begin{equation}
\hat{\boldsymbol{\lambda}}_t = H_\phi(s_t)
\end{equation}
The key to this approach is the design of the batch descriptor $s_t$, which must efficiently capture the information needed to determine the optimal expert mixture. We construct $s_t$ from two synergistic sources of information.

\paragraph{Batch Descriptor $s_t$.}
For an incoming batch $X_t$ of size $B$, we compute $s_t$ as follows:
\begin{enumerate}
    \item \textbf{Feature-level Statistics:} We pass $X_t$ through a shared, frozen feature extractor (e.g., the convolutional backbone of the expert models) to obtain a set of $B$ feature vectors, $\{f_i\}_{i=1}^B$, each of dimension $D_{feat}$. We then compute the first four statistical moments (mean, variance, skew, and kurtosis) across the batch dimension for each feature dimension. This results in a $4 \times D_{feat}$ dimensional vector that captures the low-level distributional properties of the batch in feature space. In our experiments, $D_{feat}=256$, yielding a 1024-dim vector, which we down-project to 256 dimensions using a pre-trained linear layer for compactness.
    \item \textbf{Prediction-level Statistics:} We compute predictions for the batch $X_t$ using both expert models, $\theta_A$ and $\theta_B$, yielding probability distributions $p_A(y|X_t)$ and $p_B(y|X_t)$. We also compute the entropy of the predictions for each sample. From these, we calculate statistics (mean and variance) of: (a) expert A's predictions, (b) expert B's predictions, and (c) the prediction entropy. For a $C$-class problem, this yields a $2 \times (C + C + 1) = 4C+2$ dimensional vector, capturing model confidence, disagreement, and predictive uncertainty. In our 10-class experiments, this is a 22-dim vector.
\end{enumerate}
The final batch descriptor $s_t$ is the concatenation of these two components, resulting in a $256 + 22 = 278$-dimensional vector.

\subsection{Meta-Training the Hypernetwork}
The hypernetwork $H_\phi$ is trained offline in a meta-learning phase. We generate a large meta-dataset of $(s_i, \boldsymbol{\lambda}^*_i)$ pairs.
\begin{enumerate}
    \item We create synthetic data streams by sampling batches from the \textit{training sets} of the expert domains (e.g., MNIST and FashionMNIST). This ensures zero data leakage from the final test stream.
    \item For each synthetic batch $X_i$, we compute its "oracle" merging coefficients $\boldsymbol{\lambda}^*_i$. This is done by performing a search (e.g., grid search or gradient-based optimization) to find the $\boldsymbol{\lambda}$ that minimizes the cross-entropy loss against the batch's true labels.
    \item We also compute the corresponding batch descriptor $s_i$.
    \item The hypernetwork $H_\phi$ is then trained on this meta-dataset to minimize the Mean Squared Error (MSE) between its predictions and the oracle targets:
\end{enumerate}
\begin{equation}
\label{eq:mse}
\mathcal{L}(\phi) = \frac{1}{N} \sum_{i=1}^{N} \| H_\phi(s_i) - \boldsymbol{\lambda}^*_i \|_2^2
\end{equation}
This process distills the complex, label-dependent optimization for $\boldsymbol{\lambda}^*$ into a simple, feature-driven forward pass.

\subsection{Theoretical Generalization and Loss Bound of Amortization}
In this subsection, we establish a theoretical guarantee on the performance gap of our amortized merging framework. Let $\mathcal{L}(\boldsymbol{\lambda}; X)$ denote the unsupervised test-time loss (e.g., cross-entropy or prediction entropy) evaluated on a batch $X$ with weight-merging coefficients $\boldsymbol{\lambda}$. The true optimal, label-aware oracle coefficient vector is $\boldsymbol{\lambda}^* = \arg\min_{\boldsymbol{\lambda}} \mathcal{L}(\boldsymbol{\lambda}; X)$.

Our framework replaces the iterative optimization for $\boldsymbol{\lambda}^*$ with the hypernetwork prediction $\hat{\boldsymbol{\lambda}} = H_\phi(s)$. To analyze how the prediction error of $H_\phi(s)$ affects the final model performance, we introduce the following standard Lipschitz-continuity assumption on the loss surface over the low-dimensional weight-interpolation space:

\begin{assumption}
\label{assump:lipschitz}
The test-time loss function $\mathcal{L}(\boldsymbol{\lambda}; X)$ is $L_1$-Lipschitz continuous with respect to the merging coefficients $\boldsymbol{\lambda}$ in the $\ell_2$-norm, meaning there exists a constant $L_1 > 0$ such that for any $\boldsymbol{\lambda}_1, \boldsymbol{\lambda}_2 \in [0, 1]^L$:
\begin{equation}
|\mathcal{L}(\boldsymbol{\lambda}_1; X) - \mathcal{L}(\boldsymbol{\lambda}_2; X)| \leq L_1 \|\boldsymbol{\lambda}_1 - \boldsymbol{\lambda}_2\|_2
\end{equation}
\end{assumption}

Using Assumption \ref{assump:lipschitz}, we can bound the suboptimality of the model's performance when utilizing the amortized predictions of $H_\phi$:

\begin{theorem}
Let $\epsilon > 0$ be the upper bound of the hypernetwork's prediction error on the batch descriptor $s$ under the $\ell_2$-norm, i.e., $\|H_\phi(s) - \boldsymbol{\lambda}^*\|_2 \leq \epsilon$. Then, the performance gap between the amortized Hyper-TTMM coefficients $\hat{\boldsymbol{\lambda}}$ and the label-aware oracle coefficients $\boldsymbol{\lambda}^*$ is bounded by:
\begin{equation}
\mathcal{L}(\hat{\boldsymbol{\lambda}}; X) - \mathcal{L}(\boldsymbol{\lambda}^*; X) \leq L_1 \epsilon
\end{equation}
\end{theorem}

\begin{proof}
By directly applying Assumption \ref{assump:lipschitz} with $\boldsymbol{\lambda}_1 = \hat{\boldsymbol{\lambda}} = H_\phi(s)$ and $\boldsymbol{\lambda}_2 = \boldsymbol{\lambda}^*$, we obtain:
\begin{align}
\mathcal{L}(\hat{\boldsymbol{\lambda}}; X) - \mathcal{L}(\boldsymbol{\lambda}^*; X) &\leq |\mathcal{L}(\hat{\boldsymbol{\lambda}}; X) - \mathcal{L}(\boldsymbol{\lambda}^*; X)| \\
&\leq L_1 \|\hat{\boldsymbol{\lambda}} - \boldsymbol{\lambda}^*\|_2 \\
&= L_1 \|H_\phi(s) - \boldsymbol{\lambda}^*\|_2 \\
&\leq L_1 \epsilon
\end{align}
which completes the proof.
\end{proof}

This theorem provides a solid theoretical foundation for Hyper-TTMM: it guarantees that the test-time loss of our amortized model is close to the absolute oracle minimum, and that this performance gap is directly controlled and bounded by the hypernetwork's prediction error ($\epsilon$). Consequently, minimizing the mean squared error (MSE) during our offline meta-training phase (Equation \ref{eq:mse}) directly minimizes the upper bound of the test-time loss suboptimality.

\subsection{Theoretical Computational Complexity Analysis}
The primary advantage of Hyper-TTMM is its computational efficiency. Let $F(P)$ be the FLOPs required for a single forward pass through a base model with $P$ parameters. A backward pass typically requires approximately twice the computation of a forward pass, i.e., $2 \times F(P)$.

\paragraph{Gradient-based TTA.} For an iterative method running for $K$ optimization steps, each step requires a forward pass to compute the loss (e.g., entropy) and a backward pass to compute gradients. The total complexity is:
\begin{equation}
\text{FLOPs}_{\text{TTA}} = K \times (F(P) + 2 F(P)) = 3K \cdot F(P)
\end{equation}
For a typical $K=5$ steps, this equates to roughly $15 \times F(P)$.

\paragraph{Hyper-TTMM.} Our method requires two main steps at test time: (1) a forward pass through the base model's feature extractor to compute the batch descriptor $s_t$, and (2) a forward pass through the final merged model to get predictions. The forward pass through the small hypernetwork is negligible. The complexity is:
\begin{equation}
\text{FLOPs}_{\text{Hyper-TTMM}} \approx 2 \times F(P)
\end{equation}
Comparing the two, Hyper-TTMM offers a theoretical FLOPs reduction of approximately $(15 - 2) / 15 \approx 86.7\%$. This analysis shows that by amortizing the optimization, we eliminate the expensive backward passes entirely, avoiding both the computational and memory overhead of backpropagation at test time.

\subsection{Adaptive Temporal Dynamics for Domain Shifts}
\label{sec:adaptive-ema}
To improve temporal stability and reduce noise from batch-to-batch variations, one might apply an Exponential Moving Average (EMA) to the batch descriptors: $s'_t = \alpha s'_{t-1} + (1-\alpha)s_t$. However, a fixed, high-momentum $\alpha$ (e.g., 0.9) is disastrous during an abrupt domain shift, as the smoothed state $s'_{t}$ remains "contaminated" by the previous domain, leading to highly suboptimal merging coefficients.

We propose an \textbf{Adaptive EMA with Distance-Thresholded Resetting} to overcome this. We monitor the Euclidean distance between the raw descriptors of consecutive batches:
\begin{equation}
D_t = \|s_t - s_{t-1}\|_2
\end{equation}
We observe empirically that $D_t$ is consistently small for batches within the same domain but spikes dramatically during a domain transition. We leverage this signal to dynamically control the EMA momentum $\alpha_t$. We define a threshold $\tau$ and set the update rule as:
\begin{equation}
\alpha_t =
\begin{cases}
    0 & \text{if } D_t > \tau \quad \text{(Reset on domain shift)} \\
    \alpha_{\text{smooth}} & \text{if } D_t \leq \tau \quad \text{(Smooth within domain)}
\end{cases}
\end{equation}
The smoothed descriptor is then updated as $s'_t = \alpha_t s'_{t-1} + (1 - \alpha_t)s_t$.
This allows Hyper-TTMM to react instantaneously to major distribution shifts by using the raw batch descriptor ($\alpha_t=0$), while still benefiting from noise reduction via smoothing ($\alpha_t=\alpha_{\text{smooth}}$) during stable periods. This adaptive mechanism adds negligible computational overhead.

\section{Experimental Setup}

\paragraph{Models and Datasets.} We use two expert models based on a "SimpleCNN" architecture, consisting of two convolutional layers followed by two fully connected layers. The first expert is trained on MNIST \citep{lecun1998gradient}, and the second on FashionMNIST. We use the standard training splits of these datasets to train the experts and to generate the meta-dataset for the hypernetwork.

\paragraph{Evaluation Stream.} To evaluate performance under dynamic distribution shifts, we construct a 5-phase continual test stream from the official \textit{test sets} of MNIST, FashionMNIST, and KMNIST (a novel, unseen domain). This ensures zero data leakage from the meta-training phase. The stream is ordered as follows:
\begin{itemize}
    \item \textbf{Phase 0:} Clean MNIST (in-distribution for expert A).
    \item \textbf{Phase 1:} Noisy MNIST (Gaussian noise, $\sigma=0.5$).
    \item \textbf{Phase 2:} Clean FashionMNIST (in-distribution for expert B).
    \item \textbf{Phase 3:} Noisy FashionMNIST (Gaussian noise, $\sigma=0.5$).
    \item \textbf{Phase 4:} Novel KMNIST (out-of-distribution).
\end{itemize}
This setup tests the models' ability to adapt to both domain corruption and semantic domain shifts.

\paragraph{Baselines.} We compare Hyper-TTMM against five strong baselines:
\begin{enumerate}
    \item \textbf{Expert MNIST Only:} A static model, the MNIST expert.
    \item \textbf{Expert Fashion Only:} A static model, the FashionMNIST expert.
    \item \textbf{Uniform Merging:} A static merge with $\lambda^{(l)}=0.5$ for all layers.
    \item \textbf{Gradient-based TTA:} An iterative method that performs 5 steps of gradient descent per batch to find a single global $\lambda$ that minimizes prediction entropy.
    \item \textbf{Oracle Merging (Ceiling):} A theoretical upper bound where we use true test labels to perform a grid search for the optimal layer-wise $\boldsymbol{\lambda}$ for each batch. This is computationally prohibitive but provides a performance ceiling.
\end{enumerate}

\paragraph{Hyper-TTMM Implementation.} Our default hypernetwork is a 2-layer MLP with 128 hidden units per layer and ReLU activations. It is trained on a meta-dataset of 1000 $(s_i, \boldsymbol{\lambda}^*_i)$ pairs using the Adam optimizer. For our adaptive temporal dynamics, we use $\alpha_{\text{smooth}}=0.5$ and a distance threshold of $\tau=0.3$, determined on a separate validation set. All experiments were run on a single NVIDIA H100 GPU.

\section{Experimental Results and Discussion}

\subsection{Main Performance Comparison}

Table \ref{tab:main_results} presents the main evaluation results. Hyper-TTMM achieves an outstanding overall accuracy of \textbf{50.41\%}, significantly outperforming all practical baselines and coming within a mere 0.38\% of the theoretical Oracle ceiling (50.78\%).

\begin{table*}[t]
\caption{Main evaluation results. We report phase-wise and overall accuracy (\%), and average inference latency per batch (ms). Hyper-TTMM achieves near-oracle accuracy with a tiny fraction of the latency, and massively outperforms standard TTA.}
\label{tab:main_results}
\vskip 0.15in
\begin{center}
\begin{small}
\begin{sc}
\setlength{\tabcolsep}{3.5pt}
\begin{tabular}{lccccccr}
\toprule
Method & Phase 0 & Phase 1 & Phase 2 & Phase 3 & Phase 4 & Overall & Latency \\
& (MNIST) & (Noisy M.) & (Fashion) & (Noisy F.) & (KMNIST) & Acc. (\%) & (ms) \\
\midrule
Expert MNIST & 97.03 & 30.94 & 5.78 & 10.00 & 5.62 & 29.88 & 2.84 \\
Expert Fashion & 17.03 & 10.16 & 89.38 & 29.38 & 5.94 & 30.38 & 2.83 \\
Uniform Merging & 69.22 & 19.53 & 77.03 & 13.75 & 4.84 & 36.88 & 2.82 \\
Gradient-based TTA & 75.31 & 18.44 & 78.91 & 13.91 & 4.84 & 38.28 & 45.24 \\
\midrule
Oracle Merging & 97.03 & 42.66 & 87.81 & 21.09 & 5.31 & \textbf{50.78} & 256.42 \\
\textbf{Hyper-TTMM (Ours)} & 95.00 & 42.34 & 84.84 & 24.22 & 5.62 & \underline{50.41} & \textbf{5.51} \\
\bottomrule
\end{tabular}
\end{sc}
\end{small}
\end{center}
\vskip -0.1in
\end{table*}

\paragraph{Accuracy-Latency Trade-off.}
Hyper-TTMM is \textbf{46.6x faster} than the Oracle (5.51ms vs 256.42ms) and \textbf{8.2x faster} than Gradient-based TTA (5.51ms vs 45.24ms). This confirms that amortization provides a dramatic speedup, enabling real-time dynamic model merging. The slight latency increase over static models (5.51ms vs ~2.8ms) is due to the extra forward pass for feature extraction and the hypernetwork pass, which is a small and acceptable price for a +20\% absolute accuracy gain.

\paragraph{Superiority over Iterative TTA.}
Hyper-TTMM outperforms Gradient-based TTA by \textbf{+12.13\%} absolute accuracy (50.41\% vs 38.28\%). TTA collapses on noisy phases (1 and 3), falling below Uniform Merging. This illustrates the "feedback trap", where minimizing entropy on noisy, unlabeled data causes a model to become confidently wrong. Conversely, our hypernetwork is trained offline using oracle targets. It learns a robust, non-linear mapping from batch statistics to merging coefficients, bypassing the failures of unsupervised online optimization.

\paragraph{Robustness to Noise.}
On Noisy MNIST (Phase 1), Hyper-TTMM achieves \textbf{42.34\%} accuracy, outperforming both TTA (18.44\%) and the dedicated MNIST expert (30.94\%). This demonstrates that by finding an optimal weight manifold interpolation, Hyper-TTMM creates a merged model more robust to corruption than either parent model.

\begin{table*}[t]
\caption{Ablation studies on Hypernetwork architecture capacity, input batch statistics composition, and training meta-dataset size. Validation MSE is evaluated on the validation split, and Evaluation Accuracy (\%) is measured on the 50-batch test stream.}
\label{tab:ablations_combined}
\vskip 0.15in
\begin{center}
\begin{small}
\begin{sc}
\begin{tabular}{llcc}
\toprule
Category & Configuration & Val MSE & Eval Acc. (\%) \\
\midrule
\multirow{4}{*}{Architecture} & Linear & 0.04646 & 40.06 \\
 & MLP-Small (1x64) & 0.03705 & 43.75 \\
 & MLP-Medium (2x128, Default) & \textbf{0.02223} & \underline{50.41} \\
 & MLP-Large (2x256) & 0.01966 & \textbf{50.56} \\
\midrule
\multirow{3}{*}{Input Statistics} & Prob-Ent-Only (22-dim) & 0.05143 & 37.78 \\
 & Feature-Only (256-dim) & 0.02575 & 48.81 \\
 & Full Statistics (278-dim) & \textbf{0.02223} & \textbf{50.41} \\
\midrule
\multirow{3}{*}{Dataset Size} & 250 & \textbf{0.01682} & \textbf{50.53} \\
 & 500 & 0.02169 & 50.44 \\
 & 1000 (Default) & 0.02223 & 50.41 \\
\bottomrule
\end{tabular}
\end{sc}
\end{small}
\end{center}
\vskip -0.1in
\end{table*}

\subsection{Ablation Studies}

We conducted a series of five ablation studies to validate our design choices. The key results for hypernetwork architecture capacity, input batch statistics composition, and training meta-dataset scaling are summarized in Table \ref{tab:ablations_combined}.

\subsubsection{Ablation on Hypernetwork Architecture}
We tested four hypernetwork architectures. Table \ref{tab:ablations_combined} shows that performance scales with capacity: a linear model achieves only 40.06\% accuracy, while MLP models perform significantly better. Our default MLP-Medium (2x128) provides an optimal trade-off, as the larger MLP (2x256) yields a marginal 0.15\% gain. This indicates that the 278-dim descriptor mapping can be captured effectively by a moderately-sized network. Validation MSE correlates strongly with final accuracy.

\subsubsection{Ablation on Input Batch Statistics}
We analyzed our 278-dim batch descriptor $s_t$ components. Table \ref{tab:ablations_combined} shows that both feature-level and prediction-level statistics are crucial. Prediction/entropy stats alone perform poorly (37.78\%) as they lack discriminative power. Feature representation stats alone are better (48.81\%), but combining both yields 50.41\%. This synergy shows that low-level features and high-level confidence signals provide the most comprehensive information for coefficient prediction.

\subsubsection{Ablation on Meta-Dataset Size}
We evaluated meta-training sample efficiency. Table \ref{tab:ablations_combined} shows that Hyper-TTMM is highly data-efficient: with only 250 oracle samples, it achieves 50.53\% accuracy, remaining stable up to 1000 samples. This suggests that the mapping function is smooth and learnable from very few examples, ensuring practical meta-training.

\subsubsection{Ablation on Temporal Dynamics}
Table \ref{tab:ablation_ema} highlights the importance of our adaptive temporal mechanism. While No EMA serves as our baseline, naive EMA ($\alpha=0.9$) collapses on Phase 2 (falling from 84.84\% to 54.84\%) due to transition lag. Our Adaptive EMA solves this by detecting the large descriptor distance at the domain boundary, resetting its state ($\alpha=0$), and achieving \textbf{85.00\%} on Phase 2. It also improves Noisy Fashion (Phase 3) via intra-domain smoothing, demonstrating its necessity for robust continual adaptation.

\begin{table}[h]
\caption{Ablation on temporal dynamics. Naive EMA suffers a catastrophic collapse in Phase 2 due to lag. Our Adaptive EMA solves this completely.}
\label{tab:ablation_ema}
\vskip 0.15in
\begin{center}
\begin{small}
\begin{sc}
\setlength{\tabcolsep}{2.5pt}
\begin{tabular}{lcccc}
\toprule
Method & Phase 1 & Phase 2 & Phase 3 & Overall \\
& Acc (\%) & Acc (\%) & Acc (\%) & Acc (\%) \\
\midrule
No EMA & 42.34 & \underline{84.84} & 24.22 & \textbf{50.41} \\
Naive EMA ($\alpha$=.9) & \textbf{42.50} & \textit{54.84} & 23.91 & 47.97 \\
Adaptive EMA & 42.19 & \textbf{85.00} & \textbf{25.31} & \underline{50.38} \\
\bottomrule
\end{tabular}
\end{sc}
\end{small}
\end{center}
\vskip -0.1in
\end{table}


\subsubsection{Ablation on OOD Noise Robustness}
We swept the Gaussian noise standard deviation ($\sigma$) over the test stream to evaluate OOD robustness. As shown in Table \ref{tab:ablation_noise}, Hyper-TTMM maintains a massive performance advantage across all noise levels. At $\sigma=0.3$, Hyper-TTMM achieves \textbf{73.44\%} accuracy, whereas Gradient-based TTA collapses to 33.05\%---a \textbf{40.39\%} absolute gap. This confirms that amortizing optimization via our offline oracle-trained hypernetwork is vastly more robust than relying on fragile, unsupervised entropy minimization at test time, successfully bypassing the feedback trap.

\begin{table}[hbtp]
\caption{Ablation on Out-of-Distribution (OOD) noise robustness. Accuracy (\%) is reported for varying noise standard deviation ($\sigma$) applied to the entire test stream. Hyper-TTMM demonstrates vastly superior robustness compared to all baselines.}
\label{tab:ablation_noise}
\vskip 0.15in
\begin{center}
\begin{small}
\begin{sc}
\setlength{\tabcolsep}{1.8pt}
\begin{tabular}{lccccc}
\toprule
\multirow{2}{*}{Method} & \multicolumn{5}{c}{Noise Level ($\sigma$)} \\
\cmidrule(lr){2-6}
& 0.0 & 0.3 & 0.6 & 0.9 & 1.2 \\
\midrule
Expert MNIST & 51.41 & 46.48 & 20.39 & 15.70 & 13.20 \\
Expert Fashion & 53.20 & 35.08 & 17.97 & 13.05 & 11.72 \\
Uniform Merging & 73.12 & 31.72 & 15.16 & 12.66 & 12.11 \\
Gradient-based TTA & 77.11 & 33.05 & 14.61 & 12.66 & 12.34 \\
\midrule
\textbf{Hyper-TTMM (Ours)} & \textbf{89.92} & \textbf{73.44} & \textbf{30.78} & \textbf{20.23} & \textbf{15.86} \\
\bottomrule
\end{tabular}
\end{sc}
\end{small}
\end{center}
\vskip -0.1in
\end{table}

\subsection{Evaluation on a Continuously Fluctuating Stream (CFS)}
\label{sec:cfs-eval}
To evaluate our method in a more realistic and challenging scenario, we introduce the \textbf{Continuously Fluctuating Stream (CFS)} benchmark. Real-world non-stationary streams rarely exhibit clean, block-wise boundaries. Instead, domain compositions and noise levels fluctuate continuously. We construct a 50-batch test stream where:
\begin{itemize}
    \item \textbf{Batches 0--24 (Gradual MNIST-heavy shift):} The MNIST proportion $p_t = 0.5 + 0.4 \cos(\pi t / 12.0)$ and noise standard deviation $\sigma_t = 0.3 (1.0 - \cos(\pi t / 12.0))$ fluctuate continuously.
    \item \textbf{Batches 25--39 (Sudden shift to Fashion-heavy):} At $t=25$, the MNIST proportion drops abruptly to 10\% ($p_{25}=0.1$), and then gradual shifts occur as $p_t = 0.5 - 0.4 \cos(\pi (t-25) / 7.5)$, with noise $\sigma_t = 0.3 (1.0 + \cos(\pi (t-25) / 7.5))$.
    \item \textbf{Batches 40--49 (Sudden shift to OOD KMNIST):} An abrupt domain reset occurs where the stream transitions to 100\% completely unseen KMNIST.
\end{itemize}
This setup evaluates a model's capacity to handle simultaneous continuous proportion and noise oscillations, sudden domain jumps, and novel OOD shifts.

Table \ref{tab:cfs_results} presents the results. Hyper-TTMM shows exceptional robustness. Our zero-shot method outperforms standard Gradient-based TTA by a massive \textbf{+7.16\% absolute accuracy} (41.47\% vs 34.31\%) and is \textbf{6.5x faster} (34.16ms vs 221.98ms). This confirms that our hypernetwork can generalize zero-shot to continuously oscillating mixtures of domains and noise levels that degrade online entropy-minimization.

Furthermore, introducing naive EMA smoothing ($\alpha=0.9$) drops performance to 39.62\% (a -1.85\% drop compared to no EMA) due to the moving average lag at the abrupt transition boundaries (such as $t=25$ and $t=40$). Conversely, our proposed \textbf{Adaptive EMA with Distance-Thresholded Resetting} ($\alpha=0.5$) dynamically resets its smoothing state during sudden transitions. This avoids moving average lag entirely while successfully smoothing out gradual intra-phase oscillations, resulting in the best overall practical accuracy of \textbf{41.53\%} (outperforming all non-oracle methods) with a highly optimized CPU inference latency of \textbf{27.62ms}.

\begin{table}[hbtp]
\caption{Evaluation results on our Continuously Fluctuating Stream (CFS) benchmark. We report phase accuracies (\%), overall accuracy (\%), and average CPU latency per batch (ms). MNIST, Fashion, and KMNIST refer to the three respective phases of the stream.}
\label{tab:cfs_results}
\vskip 0.15in
\begin{center}
\begin{scriptsize}
\begin{sc}
\setlength{\tabcolsep}{1.0pt}
\begin{tabular}{lccccr}
\toprule
Method & MNIST & Fashion & KMNIST & Over. (\%) & Lat. (ms) \\
\midrule
Exp. MNIST & 49.88 & 47.08 & 3.75 & 39.81 & 15.77 \\
Exp. Fashion & 29.69 & 30.42 & 5.62 & 25.09 & 16.17 \\
Uniform Merge & 42.06 & 38.75 & 5.00 & 33.66 & 16.25 \\
Grad. TTA & 42.94 & 39.58 & 4.84 & 34.31 & 221.98 \\
\midrule
Oracle Merge & 57.00 & 52.92 & 4.84 & \textbf{45.34} & 1191.66 \\
\textbf{Hyper-TTMM} & 52.06 & 49.27 & 3.28 & 41.47 & 34.16 \\
~~+ Naive EMA & 49.50 & 47.40 & 3.28 & 39.62 & 30.65 \\
\textbf{~~+ Adapt. EMA} & \textbf{52.19} & \textbf{49.17} & \textbf{3.44} & \underline{41.53} & \textbf{27.62} \\
\bottomrule
\end{tabular}
\end{sc}
\end{scriptsize}
\end{center}
\vskip -0.1in
\end{table}

\section{Conclusion, Limitations, and Future Work}

We introduced Hyper-TTMM, a novel framework that amortizes the expensive online search for model merging coefficients into a single forward pass of a lightweight hypernetwork. Hyper-TTMM achieves near-oracle performance, outperforming standard gradient-based TTA by avoiding the unsupervised "feedback trap" under noise, while remaining orders of magnitude faster. We also proposed an adaptive temporal smoothing mechanism that robustly handles abrupt domain shifts without transition lag.

\paragraph{Limitations.} Despite its strong performance and speed, Hyper-TTMM has certain limitations. First, training the hypernetwork requires an offline meta-learning phase with simulated streams, which may not capture all possible real-world out-of-distribution shifts. Second, our linear model merging assumption, although highly effective, does not handle non-linear parameter alignment, which could be beneficial when experts diverge significantly. Finally, while our Adaptive EMA successfully detects abrupt shifts on the Continuously Fluctuating Stream (CFS), selecting the optimal distance threshold $\tau$ depends on the batch statistics of the features and might require calibration for new architectures.

\paragraph{Future Work.} Future work includes scaling Hyper-TTMM to merge more than two experts, exploring non-linear merging rules, and applying this amortization principle to large-scale Vision Transformers on complex, real-world data streams. Hyper-TTMM represents a significant step towards practical, robust, and real-time adaptive model composition.

% In the unusual situation where you want a paper to appear in the
% references without citing it in the main text, use \nocite
\nocite{*}

\bibliography{icml2026_hyperttmm}
\bibliographystyle{icml2026}

%%%%%%%%%%%%%%%%%%%%%%%%%%%%%%%%%%%%%%%%%%%%%%%%%%%%%%%%%%%%%%%%%%%%%%%%%%%%%%%
%%%%%%%%%%%%%%%%%%%%%%%%%%%%%%%%%%%%%%%%%%%%%%%%%%%%%%%%%%%%%%%%%%%%%%%%%%%%%%%
% APPENDIX
%%%%%%%%%%%%%%%%%%%%%%%%%%%%%%%%%%%%%%%%%%%%%%%%%%%%%%%%%%%%%%%%%%%%%%%%%%%%%%%
%%%%%%%%%%%%%%%%%%%%%%%%%%%%%%%%%%%%%%%%%%%%%%%%%%%%%%%%%%%%%%%%%%%%%%%%%%%%%%%
% \newpage
% \appendix
% \onecolumn
% \section{You \emph{can} have an appendix here.}

% You can have as much text here as you want. The main body must be at most eight pages.
% For the final version, nine pages is allowed.
%%%%%%%%%%%%%%%%%%%%%%%%%%%%%%%%%%%%%%%%%%%%%%%%%%%%%%%%%%%%%%%%%%%%%%%%%%%%%%%
%%%%%%%%%%%%%%%%%%%%%%%%%%%%%%%%%%%%%%%%%%%%%%%%%%%%%%%%%%%%%%%%%%%%%%%%%%%%%%%


\end{document}